\section{Discussion}\label{sec:discussion}

\subsection{Methodological warnings}
Three apparent findings in this project dissolved under controlled re-analysis, each illustrating how benchmark composition can manufacture false conclusions. First, an apparent per-allele superiority over a pan model was a 10.6$\times$ training-data-size artifact, not biology. Second, a claimed ``label quality dominates architecture'' result collapsed once the degenerate random-label null (which performed at chance, $\sim$33\%) was removed --- the realistic label effect is smaller than the architecture range. Third, a reported ``simple architecture beats complex ones'' ranking was an artifact of stale model checkpoints: re-training with fixed seeds showed a CNN instead modestly outperforms the Deep FFN. We recommend that peptide-MHC benchmarks report multi-seed estimates with confidence intervals, protein-background negatives, frozen data snapshots, and a state-of-the-art comparator such as NetMHCpan, to avoid these failure modes.

\textbf{Box 1. Benchmark hygiene recommendations.} The three artifacts translate into reusable practices for immunoinformatics and machine-learning benchmarks: (1) equalise training-data size across model arms and report per-allele sample counts; (2) use a true i.i.d.\ random-label null (chance $\approx$ 33\% for three balanced classes) rather than a degenerate control; (3) freeze data snapshots, model checkpoints, and label versions, and report multi-seed mean $\pm$ SD with confidence intervals; (4) report label provenance and break label circularity with a temporal holdout and an independent state-of-the-art comparator (NetMHCpan). Each failure mode reproduced a false conclusion that was caught only by controlled re-analysis.

Architecture choice has a modest but measurable effect on peptide-MHC binding prediction. Across five seeds, a CNN (91.9\%) outperformed a Deep FFN (90.6\%) by $\sim$1.3 percentage points, while a Random Forest was clearly worse (79.7\%), for a maximum inter-architecture range of $\sim$12 percentage points. The top four architectures fell within $\sim$2 percentage points, indicating that data quality and the decision boundary matter more than architectural complexity for this task. Random-label controls performed at chance ($\sim$33\%), confirming that the observed performance reflects genuine signal rather than noise overfitting.

A random-label control performed at chance ($\sim$33\%), confirming the models do not overfit to noise, and a PSSM-derived label source yielded lower accuracy than MHCflurry labels (82.6\% vs 90.6\%, though on a separate peptide pool and therefore not a clean comparison). This does not support the hypothesis that deterministic biophysical labels are inherently superior to ML-derived labels; rather, it suggests that MHCflurry's labels are at least as learnable as a simple PSSM. We nevertheless emphasise that training on any computational label --- including MHCflurry --- re-deposits predictor bias into downstream models, and recommend that benchmarks report label provenance and the WB-class F1 as a standard diagnostic.

An independent state-of-the-art predictor, NetMHCpan-4.1, agreed with the MHCflurry labels on only 62.24\% (323/519) of the same held-out test peptides --- substantially below the 90.6\% achieved by a model trained directly on those labels. This $\sim$28-point gap quantifies the label-circularity limitation: our trained models learn MHCflurry's specific decision boundary, whereas an independent predictor disagrees on roughly one peptide in three, underscoring that the reported accuracies measure fidelity to a particular labeler rather than true binding.

BLOSUM62 proved more robust than ESM-2 at single-allele training scales; high-dimensional per-position embeddings overfit small datasets. The 2,880-dimensional per-position ESM-2 features create an unfavourable feature-to-sample ratio for the 4,671-sample training set, and the single-split 93.3\% result reflects a favourable partition rather than a genuine encoding advantage. The shallower 6-layer ESM-2 variant outperformed the deeper 12-layer model, possibly because t12's additional transformer layers encode long-range protein-level interactions that introduce noise at the short 9-mer peptide scale. Architectural complexity beyond depth with batch normalisation gave diminishing returns; a simple PSSM label source yielded lower accuracy (82.6\%) than MHCflurry labels on its own pool, showing that deep learning adds value over a hand-crafted scoring matrix for this task.

The multi-allele extension shows the recipe is portable: one rebalanced pan model matches per-allele specialists on four of five alleles (B*07:02 excepted) and reaches 0.819 pooled sensitivity on 15 alleles, approximating NetMHCpan's practical value with a simple DFFN. Yet non-canonical (motif-failing) binders plateau at $\sim$0.65 sensitivity across three independent sequence-based predictors: our Deep FFN (0.607), MHCflurry (0.668), and the ESM-2-based MUNIS (0.647) fall within a narrow band (0.61--0.67), far below the $\sim$0.98 canonical sensitivity, indicating this ceiling is shared across sequence architectures of varying complexity rather than specific to our simple network. The energy-compensated tail (P2-Lys/Arg/His reaching B-pocket Glu63) is prediction-resistant to all three. The three-tier screening pipeline we considered (sequence filtering $\rightarrow$ structural scoring $\rightarrow$ energy-minimisation validation) remains a hypothesis for these configurations, pending experimental validation, since the structural analysis presented here (Section~\ref{sec:results}) is a single-case docking hypothesis rather than a scaled benchmark. The canonical binder problem is effectively solved; the non-canonical tail defines the current frontier.

\subsection{Translational implications}
The cancer neoepitope candidates identified here are of translational interest as testable hypotheses. p53 R248W, which created MNWRPILTI, is among the most frequent TP53 missense mutations across human cancers ($\sim$3.5\% of somatic TP53 mutations, IARC) with documented gain-of-function oncogenic activity; it is among the most frequent TP53 mutations across multiple tumour types (IARC TP53 database). A vaccine targeting this neoepitope could, only if immunogenicity is experimentally confirmed, apply across multiple tumour types; validation (e.g., IFN-$\gamma$ ELISpot using HLA-A*02:01+ donor PBMCs) is required before any clinical development. KRAS G12C, which enhanced an existing epitope (WB$\rightarrow$SB), is the druggable mutation targeted by sotorasib and adagrasib, suggesting a rational combination immunotherapy strategy. The three epitope-destroying mutations (KRAS G13D, KRAS A146T, p53 R249S) represent potential immune-evasion mechanisms. We note that binding prediction alone does not establish immunogenicity: prior work shows that only a small fraction of top-predicted MHC-I binders are confirmed immunogenic, and that peptide--MHC stability predicts immunogenicity better than affinity \citep{Harndahl12,Backert15}, and the binding-to-immunogenicity gap remains the central bottleneck in translating computational epitope prediction to clinical vaccine design.

\subsection{Limitations}
Training labels are MHCflurry-derived rather than experimental affinities; we quantified the circularity by cross-referencing the 49 IEDB epitopes against MHCflurry's training corpus --- 44/49 (89.8\%) overlapped, though all five non-overlapping epitopes were correctly classified (5/5, 100\%), providing limited evidence of generalisation beyond the training distribution. The single-allele design was deliberate --- holding the allele constant isolates label-source variation from allele-driven confounds --- but the multi-allele extension addresses the generalisability question directly. The multi-allele estimates rest on IEDB-derived positives; experimental (non-IEDB) validation of the extreme tail is outstanding, and a pMHC-fine-tuned structure model (beyond generic AlphaFold-Multimer) might lift the Type B plateau. The K-E63 salt bridge, all translational predictions, and the cancer neoepitope candidates require experimental confirmation (e.g., IFN-$\gamma$ ELISpot or tetramer staining). Antigen-processing steps (proteasomal cleavage, TAP transport) are not modelled.
